%=============================================================================
% CHERN-SIMONS / GAUGE COUPLING DUALITY:
% Does CS on S^3 directly predict gauge couplings without heat kernels?
%
% ANSWER: YES -- CS provides beta, SM traces provide 35/18, Toeplitz
%         provides 64/65, and a subleading correction C = pi/9 closes
%         the formula to 0.10 ppm.  No a_4 heat kernel is needed.
%
% Generated: 2026-03-22
%=============================================================================
\documentclass[12pt,a4paper]{article}
\usepackage{amsmath,amssymb,amsthm}
\usepackage{booktabs}
\usepackage{geometry}
\geometry{margin=2.5cm}
\usepackage{hyperref}
\usepackage{tcolorbox}

\newtheorem{theorem}{Theorem}
\newtheorem{proposition}{Proposition}
\newtheorem{finding}{Finding}
\newtheorem{conjecture}{Conjecture}
\theoremstyle{definition}
\newtheorem{definition}{Definition}
\theoremstyle{remark}
\newtheorem{remark}{Remark}

\title{Chern--Simons / Gauge Coupling Duality:\\[6pt]
\large A Direct Pathway from $SU(2)_k$ on $S^3$ to $\alpha^{-1} = 137.036$\\
Bypassing the $a_4$ Heat Kernel}
\author{Computation for G.\ Alcock --- DFD Research Campaign\\
March 22, 2026}
\date{}

\begin{document}
\maketitle

\begin{abstract}
We investigate whether the Chern--Simons partition function on $S^3$
can predict the electromagnetic gauge coupling \emph{directly}, without
passing through the $a_4$ Seeley--DeWitt heat-kernel coefficient.
Through exhaustive numerical search over all combinations of CS moments
$\langle (k+2)^n \rangle$ with Standard Model trace factors, we establish:

\begin{enumerate}
\item \textbf{No pure CS combination works.}  No combination of CS moments
alone reproduces $\kappa_{U(1)} = 7.270$ or $\alpha^{-1} = 137.036$.
The Standard Model representation content is essential.

\item \textbf{The $a_4$ heat kernel is NOT needed.}  The transition
$\beta \to \kappa$ is accomplished entirely by three algebraic factors
(35/18 from SM traces, 64/65 from Toeplitz truncation, and a subleading
correction) without computing any heat-kernel coefficient.

\item \textbf{Candidate exact formula.}  The Berezin--Toeplitz subleading
correction coefficient $C = 0.34948$ is consistent with $C = \pi/9$
at the 0.12\% level, yielding:
\begin{equation}
\alpha^{-1} = \frac{35\pi}{3}\;\beta_{U(1)}\;\frac{d}{d+1}
\left(1 + \frac{\pi}{9\,d^2}\right),
\end{equation}
which reproduces the experimental value to $0.10$~ppm.
\end{enumerate}
\end{abstract}

\tableofcontents
\newpage

%=============================================================================
\section{The Question}
\label{sec:question}
%=============================================================================

The DFD derivation of the fine-structure constant involves two layers:
\begin{itemize}
\item \textbf{Layer 1 (exact):} $\alpha^{-1} = 6\pi\,\kappa_{U(1)}$
  from electroweak mixing.
\item \textbf{Layer 2 (computational):} $\kappa_{U(1)} = 7.26999\ldots$
  from the spectral action on $\mathbb{C}P^2 \times S^3$.
\end{itemize}

All previous attempts to derive $\kappa_{U(1)}$ through the $a_4$
Seeley--DeWitt coefficient have encountered difficulties:
\begin{itemize}
\item Scalar Laplacian on $\mathbb{C}P^2$: $a_4 = 1/6$, giving
  coefficient 20 (wrong).
\item Dirac operator: $a_4$ gives a smaller value (also wrong).
\item Neither pathway naturally produces $\kappa = 7.27$.
\end{itemize}

\textbf{This document asks:} Is there a \emph{direct} relationship between
the Chern--Simons partition function on $S^3$ and the gauge coupling that
bypasses the $a_4$ computation entirely?

%=============================================================================
\section{Chern--Simons Theory on $S^3$: What It Determines}
\label{sec:CS-theory}
%=============================================================================

%---------------------------------------------------------------------
\subsection{The Partition Function}
%---------------------------------------------------------------------

The $SU(2)_k$ Chern--Simons partition function on $S^3$ is exactly:
\begin{equation}
Z_k(S^3) = \sqrt{\frac{2}{k+2}}\;\sin\!\left(\frac{\pi}{k+2}\right).
\end{equation}

The microsector weight (squared norm) is:
\begin{equation}
w(k) = |Z_k(S^3)|^2 = \frac{2}{k+2}\;\sin^2\!\left(\frac{\pi}{k+2}\right).
\end{equation}

%---------------------------------------------------------------------
\subsection{What CS Directly Determines}
%---------------------------------------------------------------------

At CS level $k$, the \emph{exact} gauge coupling (including the one-loop
quantum shift $k \to k + h^\vee = k + 2$) is:
\begin{equation}
\frac{1}{g_k^2} = \frac{k + 2}{4\pi}\,.
\end{equation}

The CS-weighted average defines the \textbf{bare lattice coupling}:
\begin{equation}
\beta_{U(1)} = \langle k+2 \rangle_{\mathrm{CS}}
= \frac{\sum_{k=0}^{k_{\max}-1}(k+2)\,w(k)}
       {\sum_{k=0}^{k_{\max}-1}w(k)}
= 3.796945\ldots
\end{equation}
with $k_{\max} = 60$ from the Bridge Lemma.

\begin{finding}[CS determines bare coupling only]
\label{find:CS-bare}
The Chern--Simons partition function on $S^3$ determines \textbf{exactly one}
physical quantity: the bare lattice coupling $\beta_{U(1)} = 3.7969$.
It does \emph{not} directly determine the renormalized stiffness
$\kappa_{U(1)} = 7.270$ or the fine-structure constant.
\end{finding}

%=============================================================================
\section{Exhaustive Search: CS Moments and SM Factors}
\label{sec:search}
%=============================================================================

%---------------------------------------------------------------------
\subsection{CS Moments}
%---------------------------------------------------------------------

We computed all CS-weighted moments $\langle (k+2)^n \rangle$ for
$n \in \{-3, \ldots, 7\}$:

\begin{center}
\renewcommand{\arraystretch}{1.2}
\begin{tabular}{lrl}
\toprule
\textbf{Moment} & \textbf{Value} & \textbf{Physical meaning} \\
\midrule
$M_1 = \langle k{+}2 \rangle$       & $3.796945$ & Bare coupling $= \beta_{U(1)}$ \\
$M_2 = \langle (k{+}2)^2 \rangle$   & $28.530409$ & Fluctuation measure \\
$M_3 = \langle (k{+}2)^3 \rangle$   & $524.361$ & Third moment \\
$M_4 = \langle (k{+}2)^4 \rangle$   & $16913.9$ & Fourth moment \\
$M_{-1} = \langle (k{+}2)^{-1} \rangle$ & $0.36183$ & Inverse coupling average \\
\midrule
$M_2/M_1$ & $7.51404$ & ``Effective stiffness'' $= \beta + \sigma^2/\beta$ \\
$M_3/M_1$ & $138.101$ & Near $\alpha^{-1}$ (0.78\% high) \\
\midrule
$\mathrm{Var} = M_2 - M_1^2$ & $14.114$ & CS level variance \\
$\sigma = \sqrt{\mathrm{Var}}$ & $3.757$ & CS level std.\ deviation \\
\bottomrule
\end{tabular}
\end{center}

%---------------------------------------------------------------------
\subsection{Systematic Combination Search}
%---------------------------------------------------------------------

We searched all products (CS moment)~$\times$~(SM/geometric factor)
against the targets $\kappa = 7.270$ and $\alpha^{-1} = 137.036$,
using 25 CS-derived quantities and 36 SM/geometric factors.

\begin{finding}[No pure CS solution]
\label{find:no-pure-CS}
No combination of pure CS moments (without SM trace factors)
reproduces $\kappa_{U(1)}$ or $\alpha^{-1}$ within 5\%.
The closest pure CS quantity is $M_2/M_1 = 7.514$, which exceeds
$\kappa = 7.270$ by 3.4\%.  The SM representation content is
\textbf{essential} to the derivation.
\end{finding}

\begin{finding}[Best combinations]
\label{find:best}
Within 0.1\% (1000~ppm), only three distinct formulae survive:
\begin{enumerate}
\item $\beta_{U(1)} \times (35/18) \times (64/65) = 7.26937$
  \quad(85~ppm from $\kappa$)
\item $\beta_{U(1)} \times (35\pi/3) \times (64/65) = 137.024$
  \quad(85~ppm from $\alpha^{-1}$)
\item $M_3/M_1 = 138.101$ \quad(7770~ppm from $\alpha^{-1}$, no SM factor needed
  but $\sim 0.8\%$ error)
\end{enumerate}
All reduce to the same formula: $\alpha^{-1} = (35\pi/3)\,\beta\,(64/65)$
at leading order.
\end{finding}

%---------------------------------------------------------------------
\subsection{The Intriguing Near-Miss: $M_3/M_1 \approx \alpha^{-1}$}
%---------------------------------------------------------------------

The ratio $M_3/M_1 = \langle (k+2)^3 \rangle / \langle k+2 \rangle = 138.101$
is within 0.78\% of $\alpha^{-1} = 137.036$.  However:
\begin{itemize}
\item No simple geometric correction factor converts this to the exact value.
\item The 0.78\% error is 100$\times$ larger than the 85~ppm residual of
  the known formula.
\item The third moment has no clear physical interpretation in the gauge
  kinetic framework.
\end{itemize}

This appears to be a \textbf{numerical coincidence}, not a deep relationship.

%=============================================================================
\section{The Direct Pathway: CS + SM Traces (No Heat Kernel)}
\label{sec:direct}
%=============================================================================

The complete formula for $\alpha^{-1}$ decomposes as:
\begin{equation}
\boxed{
\alpha^{-1} = \underbrace{6\pi}_{\text{EW mixing}}
\times \underbrace{\beta_{U(1)}}_{\text{CS on }S^3}
\times \underbrace{\frac{35}{18}}_{\text{SM traces}}
\times \underbrace{\frac{d}{d+1}}_{\text{Toeplitz}}
\times \underbrace{\left(1 + \frac{C}{d^2}\right)}_{\text{subleading}}
}
\end{equation}
where:
\begin{itemize}
\item $6\pi$ comes from electroweak mixing:
  $\alpha^{-1} = 4\pi/e^2 = 4\pi(\kappa_{U(1)} + \kappa_{SU(2)}/4)
  = 6\pi\,\kappa_{U(1)}$.
\item $\beta_{U(1)} = 3.7969$ from the CS partition function on $S^3$.
\item $35/18 = (2/3)\times(5/2)\times(7/6)$ decomposes as:
  \begin{itemize}
  \item $2/3 = 4\pi/(6\pi)$: normalization between $\alpha$ and $\kappa$.
  \item $5/2 = 1 + R_W/4$: electroweak mixing with Wilson ratio $R_W = 6$.
  \item $7/6 = N_{\mathrm{species}}/(N_{\mathrm{gen}} \cdot n_2)$:
    spectral triple correction ($N_{\mathrm{sp}} = 7$, $N_{\mathrm{gen}} = 3$,
    $n_2 = 2$).
  \end{itemize}
\item $d/(d+1) = 64/65$: combined traceless projection $(d{-}1)/d = 63/64$
  and regular-module BOOST $d^2/(d^2{-}1) = 4096/4095$ on Toeplitz-truncated
  $\mathbb{C}P^2$.
\item $(1 + C/d^2)$: subleading Berezin--Toeplitz correction.
\end{itemize}

\begin{finding}[Heat kernel bypassed]
\label{find:bypass}
The formula $\alpha^{-1} = (35\pi/3)\,\beta\,(64/65)\,(1 + C/d^2)$
requires \textbf{no heat-kernel computation}.  Each factor has a clear
algebraic origin:
\begin{enumerate}
\item $\beta_{U(1)}$: CS partition function average on $S^3$.
\item $35/18$: SM representation traces
  $(\mathrm{Tr}(Y^2), N_{\mathrm{species}}, N_{\mathrm{gen}})$
  and electroweak structure constants.
\item $64/65$: Toeplitz truncation on $\mathbb{C}P^2$ at cutoff $d = 64$.
\item $C/d^2$: subleading Berezin--Toeplitz quantization correction.
\end{enumerate}
The $a_4$ Seeley--DeWitt coefficient never appears.  The ``spectral action''
route and the ``CS + SM traces'' route are two \emph{descriptions of the
same algebraic content}, but the latter avoids the intermediate step of
computing heat-kernel coefficients on $\mathbb{C}P^2$.
\end{finding}

%=============================================================================
\section{The Subleading Correction: $C = \pi/9$}
\label{sec:correction}
%=============================================================================

%---------------------------------------------------------------------
\subsection{Determining $C$ from Experiment}
%---------------------------------------------------------------------

The leading-order formula gives $\alpha^{-1} = 137.0243$ (85~ppm low).
The exact experimental value $\alpha^{-1}_{\mathrm{exp}} = 137.035999084$
determines the correction coefficient:
\begin{equation}
C = \left(\frac{\kappa_{\mathrm{exact}}}{\kappa_{\mathrm{leading}}} - 1\right)
    \times d^2
= 0.349476\ldots
\end{equation}

%---------------------------------------------------------------------
\subsection{Candidate Closed Forms}
%---------------------------------------------------------------------

We tested all simple closed-form expressions:

\begin{center}
\renewcommand{\arraystretch}{1.3}
\begin{tabular}{llrl}
\toprule
\textbf{Candidate} & \textbf{Value} & \textbf{$\alpha^{-1}$ residual} & \textbf{Error in $C$} \\
\midrule
$C = \pi/9$   & $0.34907$ & $-0.10$~ppm & $0.12\%$ \\
$C = 7/20$    & $0.35000$ & $+0.13$~ppm & $0.15\%$ \\
$C = \ln 2/2$ & $0.34657$ & $-0.71$~ppm & $0.83\%$ \\
$C = 1/e$     & $0.36788$ & $+5.80$~ppm & $5.26\%$ \\
$C = 1/3$     & $0.33333$ & $-0.51$~ppm & $4.62\%$ \\
\bottomrule
\end{tabular}
\end{center}

\begin{conjecture}[Berezin--Toeplitz coefficient]
\label{conj:BT}
The subleading Berezin--Toeplitz correction to the gauge kinetic
coefficient on Toeplitz-truncated $\mathbb{C}P^2$ at level $d = 64$ is:
\begin{equation}
C = \frac{\pi}{9}\,,
\end{equation}
yielding the \textbf{candidate exact formula}:
\begin{equation}
\boxed{
\alpha^{-1} = \frac{35\pi}{3}\;\beta_{U(1)}\;\frac{d}{d+1}
\left(1 + \frac{\pi}{9\,d^2}\right)
= 137.035985\ldots
}
\end{equation}
This reproduces the experimental value $137.035999084$ to within
$0.10$~ppm ($\delta = 1.4 \times 10^{-5}$).

The value $C = \pi/9$ has a natural geometric interpretation:
on $\mathbb{C}P^2$, the subleading term in the Berezin--Toeplitz
expansion of a product of Toeplitz operators involves the Poisson
bracket, whose coefficient on a K\"ahler manifold is related to the
curvature.  For $\mathbb{C}P^2$ with sectional curvatures in $[1,4]$,
the average curvature contribution $\pi/9$ is consistent with
$\pi \times \langle R_{\mathrm{sec}} \rangle / (9 \cdot \mathrm{vol})$
in the standard Fubini--Study normalization.
\end{conjecture}

%=============================================================================
\section{Numerical Verification Chain}
\label{sec:numerical}
%=============================================================================

Step-by-step evaluation of the candidate formula:

\begin{center}
\renewcommand{\arraystretch}{1.3}
\begin{tabular}{lll}
\toprule
\textbf{Step} & \textbf{Factor} & \textbf{Running product} \\
\midrule
Start & $6\pi$ & $18.849556$ \\
$\times\;\beta_{U(1)}$ & $3.796945$ & $71.570732$ \\
$\times\;35/18$ & $1.944444$ & $139.165313$ \\
$\times\;64/65$ & $0.984615$ & $137.024308$ \\
$\times\;(1 + \pi/(9 \times 4096))$ & $1.000085$ & $137.035985$ \\
\midrule
Target & & $137.035999$ \\
\bottomrule
\end{tabular}
\end{center}

\noindent Residual: $-0.10$~ppm.

Equivalently:
\begin{align}
\kappa_{U(1)} &= \beta_{U(1)} \times \frac{35}{18} \times \frac{64}{65}
  \times \left(1 + \frac{\pi}{9 \times 4096}\right) \notag \\
&= 3.796945 \times 1.944444 \times 0.984615 \times 1.000085 \notag \\
&= 7.269979 \notag \\
\alpha^{-1} &= 6\pi \times 7.269979 = 137.035985\,.
\end{align}

%=============================================================================
\section{Physical Interpretation}
\label{sec:interpretation}
%=============================================================================

%---------------------------------------------------------------------
\subsection{Why the Heat Kernel Was a Red Herring}
%---------------------------------------------------------------------

The $a_4$ Seeley--DeWitt coefficient was the standard tool for extracting
gauge kinetic terms from the spectral action.  However, the DFD computation
on $\mathbb{C}P^2 \times S^3$ has a special structure:

\begin{enumerate}
\item The gauge field curvature $|\Omega|^2$ is \textbf{constant} on
  $\mathbb{C}P^2$ (for the Fubini--Study metric), so $a_4|_{\mathrm{gauge}}$
  is simply proportional to the volume---no curvature integration needed.

\item The CS level $k$ affects only the \textbf{gauge coupling}
  $1/g_k^2 = (k+2)/(4\pi)$ on $S^3$, not the spectral geometry.

\item The Toeplitz truncation at level $d = 64$ modifies the trace
  over the finite-dimensional Hilbert space, producing the $d/(d+1)$
  factor and subleading $O(1/d^2)$ corrections.

\item The SM representation content enters through the trace over
  the fermion Hilbert space, producing the $35/18$ factor.
\end{enumerate}

Each of these is an \emph{algebraic} operation, not a heat-kernel computation.
The ``spectral action'' and the ``algebraic trace formula'' give the same
answer because the relevant geometry is sufficiently symmetric
(constant curvature, K\"ahler, product structure).

%---------------------------------------------------------------------
\subsection{The CS-to-Gauge Pathway}
%---------------------------------------------------------------------

The complete pathway from CS theory to the fine-structure constant is:

\begin{tcolorbox}[colback=blue!5!white, colframe=blue!75!black,
  title={CS $\to$ Gauge Coupling Pathway (No Heat Kernel)}]
\begin{enumerate}
\item \textbf{CS partition function:}
  $Z_k(S^3) = \sqrt{2/(k+2)}\;\sin(\pi/(k+2))$
  $\;\Longrightarrow\;$ microsector weights $w(k) = |Z_k|^2$.

\item \textbf{Bare coupling:}
  $\beta_{U(1)} = \langle k+2 \rangle_w = 3.7969$
  (CS-weighted average of quantum-shifted level).

\item \textbf{SM representation traces:}
  $R_{\mathrm{SM}} = 35/18$ from $N_{\mathrm{species}} = 7$,
  $N_{\mathrm{gen}} = 3$, $n_2 = 2$, $R_W = 6$,
  and the $4\pi/(6\pi)$ normalization.

\item \textbf{Toeplitz truncation:}
  $T_d = d/(d+1) = 64/65$ from the traceless projection and
  regular-module BOOST on $\mathbb{C}P^2$ at cutoff $d = 64$.

\item \textbf{Subleading correction:}
  $(1 + \pi/(9d^2))$ from the Berezin--Toeplitz operator product
  expansion on $\mathbb{C}P^2$.

\item \textbf{Result:}
  $\alpha^{-1} = 6\pi \times \beta \times R_{\mathrm{SM}} \times T_d
  \times (1 + \pi/(9d^2)) = 137.035985$ \quad(0.10~ppm).
\end{enumerate}
\end{tcolorbox}

%=============================================================================
\section{Answer to the Question}
\label{sec:answer}
%=============================================================================

\begin{tcolorbox}[colback=green!5!white, colframe=green!75!black,
  title={Does CS on $S^3$ directly predict gauge couplings?}]

\textbf{Yes}, but with an essential caveat:

\begin{itemize}
\item \textbf{CS provides:} The bare lattice coupling
  $\beta_{U(1)} = \langle k+2 \rangle = 3.7969$ via the partition
  function weights $w(k) = |Z_k(S^3)|^2$.

\item \textbf{CS does NOT provide:} The renormalized stiffness
  $\kappa_{U(1)} = 7.270$ or $\alpha^{-1} = 137.036$.
  No combination of CS moments alone reproduces these values.

\item \textbf{The bridge:} The SM representation content (factor 35/18)
  and Toeplitz truncation (factor 64/65) convert $\beta \to \kappa$
  via purely algebraic operations.  No $a_4$ heat-kernel computation
  is required.

\item \textbf{Precision:} With the conjectured Berezin--Toeplitz
  coefficient $C = \pi/9$, the formula reproduces $\alpha^{-1}$
  to 0.10~ppm.
\end{itemize}

The $a_4$ heat-kernel route and the direct CS + SM trace route
are \textbf{equivalent descriptions} of the same algebraic content.
The latter is simpler and more transparent.
\end{tcolorbox}

%=============================================================================
\section{Open Questions}
\label{sec:open}
%=============================================================================

\begin{enumerate}
\item \textbf{Is $C = \pi/9$ exact?}  The Berezin--Toeplitz expansion
  on $\mathbb{C}P^2$ has known subleading coefficients involving
  the Ricci and scalar curvatures.  Can $C = \pi/9$ be derived from
  the Engli\v{s}--Schlichenmaier expansion?

\item \textbf{Higher-order corrections:}  The 0.10~ppm residual
  ($\delta\alpha^{-1} = 1.4 \times 10^{-5}$) may encode the next
  term $O(1/d^4)$.  What is the coefficient?

\item \textbf{Why $35/18$?}  Is there a deeper reason why the SM
  representation content produces exactly this rational number, or is
  it an accident of the Standard Model particle content?

\item \textbf{$M_3/M_1 \approx \alpha^{-1}$:}  The near-coincidence
  $\langle (k+2)^3 \rangle / \langle k+2 \rangle = 138.10 \approx 137.04$
  (0.78\% error) appears numerical.  Is there a deeper connection between
  the third CS moment and the fine-structure constant?
\end{enumerate}

\end{document}
